\documentclass{amsart}
\usepackage[UTF8,fontset=windows]{ctex}
\usepackage[a4paper,margin=24mm]{geometry}
% For dealing with references we use the comment environment
\usepackage{verbatim}
\newenvironment{reference}{\comment}{\endcomment}
%\newenvironment{reference}{}{}
\newenvironment{slogan}{\comment}{\endcomment}
\newenvironment{history}{\comment}{\endcomment}

% For commutative diagrams we use Xy-pic
\usepackage[all]{xy}

% We use 2cell for 2-commutative diagrams.
\xyoption{2cell}
\UseAllTwocells

% We use multicol for the list of chapters between chapters
\usepackage{multicol}

% This is generally recommended for better output
\usepackage{lmodern}
\usepackage[T1]{fontenc}

% For cross-file-references

% Package for hypertext links:
\usepackage{hyperref}

% For any local file, say "hello.tex" you want to link to please

% Theorem environments.
%
\theoremstyle{plain}
\newtheorem{theorem}[subsection]{Theorem}
\newtheorem{proposition}[subsection]{Proposition}
\newtheorem{lemma}[subsection]{Lemma}

\theoremstyle{definition}
\newtheorem{definition}[subsection]{Definition}
\newtheorem{example}[subsection]{Example}
\newtheorem{exercise}[subsection]{Exercise}
\newtheorem{situation}[subsection]{Situation}

\theoremstyle{remark}
\newtheorem{remark}[subsection]{Remark}
\newtheorem{remarks}[subsection]{Remarks}

\numberwithin{equation}{subsection}

% Macros
%
\def\lim{\mathop{\mathrm{lim}}\nolimits}
\def\colim{\mathop{\mathrm{colim}}\nolimits}
\def\Spec{\mathop{\mathrm{Spec}}}
\def\Hom{\mathop{\mathrm{Hom}}\nolimits}
\def\Ext{\mathop{\mathrm{Ext}}\nolimits}
\def\SheafHom{\mathop{\mathcal{H}\!\mathit{om}}\nolimits}
\def\SheafExt{\mathop{\mathcal{E}\!\mathit{xt}}\nolimits}
\def\Sch{\mathit{Sch}}
\def\Mor{\mathop{\mathrm{Mor}}\nolimits}
\def\Ob{\mathop{\mathrm{Ob}}\nolimits}
\def\Sh{\mathop{\mathit{Sh}}\nolimits}
\def\NL{\mathop{N\!L}\nolimits}
\def\CH{\mathop{\mathrm{CH}}\nolimits}
\def\proetale{{pro\text{-}\acute{e}tale}}
\def\etale{{\acute{e}tale}}
\def\QCoh{\mathit{QCoh}}
\def\Ker{\mathop{\mathrm{Ker}}}
\def\Im{\mathop{\mathrm{Im}}}
\def\Coker{\mathop{\mathrm{Coker}}}
\def\Coim{\mathop{\mathrm{Coim}}}

% Boxtimes
%
\DeclareMathSymbol{\boxtimes}{\mathbin}{AMSa}{"02}

%
% Macros for moduli stacks/spaces
%
\def\QCohstack{\mathcal{QC}\!\mathit{oh}}
\def\Cohstack{\mathcal{C}\!\mathit{oh}}
\def\Spacesstack{\mathcal{S}\!\mathit{paces}}
\def\Quotfunctor{\mathrm{Quot}}
\def\Hilbfunctor{\mathrm{Hilb}}
\def\Curvesstack{\mathcal{C}\!\mathit{urves}}
\def\Polarizedstack{\mathcal{P}\!\mathit{olarized}}
\def\Complexesstack{\mathcal{C}\!\mathit{omplexes}}
% \Pic is the operator that assigns to X its picard group, usage \Pic(X)
% \Picardstack_{X/B} denotes the Picard stack of X over B
% \Picardfunctor_{X/B} denotes the Picard functor of X over B
\def\Pic{\mathop{\mathrm{Pic}}\nolimits}
\def\Picardstack{\mathcal{P}\!\mathit{ic}}
\def\Picardfunctor{\mathrm{Pic}}
\def\Deformationcategory{\mathcal{D}\!\mathit{ef}}

\hypersetup{hidelinks}
\widowpenalty=10000\clubpenalty=10000\displaywidowpenalty=10000
\begin{document}
\title[Stacks Project, Tag 0FD8]{深度二条件下充分负扭的第一上同调消失}
\maketitle
\noindent\textbf{Editorial reading note.} Source symbols, including the printed P\textasciicircum{}1 in the hyperplane paragraph and exponents in the diagram, are preserved; editorial corrections are not inserted into author text.


\noindent\textbf{Counting scope.} The primary problem is Tag 0FD8, marked below. All preceding numbered statements are supporting author context and are not counted as additional problems. This is a local pilot with a temporary ID, not production admission.


\noindent Copyright (C) 2005--2025 Johan de Jong. Stacks Project authors. Source: \href{https://stacks.math.columbia.edu/tag/0FD8}{Tag 0FD8}.

\noindent Difficulty H2: 本题在闭点处depth至少二时证明充分负的 ample 扭曲具有零H1。证明先用超平面截面使H1维数稳定，再在高次截面的核空间中选避开伴随点的截面，令同一乘法同时为零且单射以迫使稳定维数为零。该反证构造和有限域下降均未被背景直接给出，判H2。.

\noindent History: excerpt from revision local cached snapshot; exact chapter hashes; revision unknown, varieties.tex, lines 11157--11290. Only reference presentation was adapted. Licensed under GNU FDL 1.2 or later; no invariant sections or cover texts. See accompanying COPYING.

\begin{lemma}
\label{lemma-exact-sequence-induction}
Let $k$ be an infinite field. Let $n \geq 1$.
Let $\mathcal{F}$ be a coherent module on $\mathbf{P}^n_k$.
Then there exist a nonzero section
$s \in \Gamma(\mathbf{P}^n_k, \mathcal{O}(1))$
and a short exact sequence
$$
0 \to \mathcal{F}(-1) \to \mathcal{F} \to i_*\mathcal{G} \to 0
$$
where $i : H \to \mathbf{P}^n_k$ is the hyperplane $H$ associated to $s$
and $\mathcal{G} = i^*\mathcal{F}$.
\end{lemma}

\begin{proof}
The map $\mathcal{F}(-1) \to \mathcal{F}$ comes from
Constructions, Equation (\href{https://stacks.math.columbia.edu/tag/03GJ}{equation multiply on sheaf (Tag 03GJ)})
with $n = 1$, $m = -1$ and the section $s$ of $\mathcal{O}(1)$.
Let's work out what this map looks like if we restrict it to
$D_{+}(T_0)$. Write $D_{+}(T_0) = \Spec(k[x_1, \ldots, x_n])$
with $x_i = T_i/T_0$. Identify $\mathcal{O}(1)|_{D_{+}(T_0)}$ with
$\mathcal{O}$ using the section $T_0$. Hence if
$s = \sum a_iT_i$ then $s|_{D_{+}(T_0)} = a_0 + \sum a_ix_i$
with the identification chosen above. Furthermore, suppose
$\mathcal{F}|_{D_{+}(T_0)}$ corresponds to the finite
$k[x_1, \ldots, x_n]$-module $M$. Via the identification
$\mathcal{F}(-1) = \mathcal{F} \otimes \mathcal{O}(-1)$
and our chosen trivialization of $\mathcal{O}(1)$ we see that
$\mathcal{F}(-1)$ corresponds to $M$ as well. Thus
restricting $\mathcal{F}(-1) \to \mathcal{F}$ to $D_{+}(T_0)$ gives
the map
$$
M \xrightarrow{a_0 + \sum a_ix_i} M
$$
To see that the arrow is injective, it suffices to pick
$a_0 + \sum a_ix_i$ outside any of the associated primes of $M$, see
Algebra, Lemma \href{https://stacks.math.columbia.edu/tag/00LD}{lemma ass zero divisors (Tag 00LD)}. By
Algebra, Lemma \href{https://stacks.math.columbia.edu/tag/00LC}{lemma finite ass (Tag 00LC)}
the set $\text{Ass}(M)$ of associated primes of $M$ is finite.
Note that for $\mathfrak p \in \text{Ass}(M)$ the intersection
$\mathfrak p \cap \{a_0 + \sum a_i x_i\}$ is a proper $k$-subvector space.
We conclude that there is a finite family of proper sub vector spaces
$V_1, \ldots, V_m \subset \Gamma(\mathbf{P}^n_k, \mathcal{O}(1))$
such that if we take $s$ outside of $\bigcup V_i$, then multiplication
by $s$ is injective over $D_{+}(T_0)$.
Similarly for the restriction to $D_{+}(T_j)$ for $j = 1, \ldots, n$.
Since $k$ is infinite, a finite union of proper sub vector spaces
is never equal to the whole space, hence we may choose $s$ such
that the map is injective.
The cokernel of $\mathcal{F}(-1) \to \mathcal{F}$ is annihilated
by $\Im(s : \mathcal{O}(-1) \to \mathcal{O})$ which is
the ideal sheaf of $H$ by
Divisors, Definition \ref{divisors-definition-zero-scheme-s}.
Hence we obtain $\mathcal{G}$ on $H$ using
Cohomology of Schemes, Lemma \href{https://stacks.math.columbia.edu/tag/087T}{lemma i star equivalence (Tag 087T)}.
\end{proof}


\begin{lemma}
\label{lemma-vanishin-h0-negative}
Let $k$ be a field. Let $X$ be a proper scheme over $k$. Let $\mathcal{L}$
be an ample invertible $\mathcal{O}_X$-module. Let $\mathcal{F}$ be a
coherent $\mathcal{O}_X$-module. If $\text{Ass}(\mathcal{F})$ does not
contain any closed points, then
$\Gamma(X, \mathcal{F} \otimes_{\mathcal{O}_X} \mathcal{L}^{\otimes n}) = 0$
for $n \ll 0$.
\end{lemma}

\begin{proof}
For a coherent $\mathcal{O}_X$-module $\mathcal{F}$
let $\mathcal{P}(\mathcal{F})$ be the property: there
exists an $n_0 \in \mathbf{Z}$ such that for $n \leq n_0$
every section $s$ of
$\mathcal{F} \otimes_{\mathcal{O}_X} \mathcal{L}^{\otimes n}$
has support consisting only of closed points.
Since $\text{Ass}(\mathcal{F}) =
\text{Ass}(\mathcal{F} \otimes_{\mathcal{O}_X} \mathcal{L}^{\otimes n})$
we see that it suffices to prove $\mathcal{P}$
holds for all coherent modules on $X$.
To do this we will prove that conditions (1), (2), and (3) of
Cohomology of Schemes, Lemma
\href{https://stacks.math.columbia.edu/tag/01YM}{lemma property higher rank cohomological (Tag 01YM)}
are satisfied.

\medskip\noindent
To see condition (1) suppose that
$$
0 \to \mathcal{F}_1 \to \mathcal{F} \to \mathcal{F}_2 \to 0
$$
is a short exact sequence of coherent $\mathcal{O}_X$-modules
such that we have $\mathcal{P}$ for $\mathcal{F}_i$, $i = 1, 2$.
Let $n_1, n_2$ be the cutoffs we find.
Let $\mathcal{F}'_2 \subset \mathcal{F}_2$ be the maximal
coherent submodule whose support is a finite set of closed
points. Let $\mathcal{I} \subset \mathcal{O}_X$ be the annihilator
of $\mathcal{F}'_2$. Since $\mathcal{L}$ is ample, we can find an $e > 0$
such that $\mathcal{I} \otimes_{\mathcal{O}_X} \mathcal{L}^{\otimes e}$
is globally generated. Set $n_0 = \min(n_2, n_1 - e)$. Let $n \leq n_0$
and let $t$ be a global section of
$\mathcal{F} \otimes \mathcal{L}^{\otimes n}$.
The image of $t$ in $\mathcal{F}_2 \otimes \mathcal{L}^{\otimes n}$
falls into $\mathcal{F}'_2 \otimes \mathcal{L}^{\otimes n}$ because
$n \leq n_2$. Hence for any
$s \in \Gamma(X, \mathcal{I} \otimes_{\mathcal{O}_X} \mathcal{L}^{\otimes e})$
the product $t \otimes s$ lies in
$\mathcal{F}_1 \otimes \mathcal{L}^{\otimes n + e}$.
Thus $t \otimes s$ has support contained in
the finite set of closed points in $\text{Ass}(\mathcal{F}_1)$
because $n + e \leq n_1$. Since by our choice of $e$
we may choose $s$ invertible in any point
not in the support of $\mathcal{F}'_2$
we conclude that the support of $t$ is contained in the union of the 
finite set of closed points in $\text{Ass}(\mathcal{F}_1)$ and
the finite set of closed points in $\text{Ass}(\mathcal{F}_2)$.
This finishes the proof of condition (1).

\medskip\noindent
Condition (2) is immediate.

\medskip\noindent
For condition (3) we choose $\mathcal{G} = \mathcal{O}_Z$.
In this case, if $Z$ is a closed point of $X$, then there is
nothing the show. If $\dim(Z) > 0$, then we will show that
$\Gamma(Z, \mathcal{L}^{\otimes n}|_Z) = 0$
for $n < 0$. Namely, let $s$ be a nonzero section of a negative
power of $\mathcal{L}|_Z$. Choose a nonzero section $t$ of a positive
power of $\mathcal{L}|_Z$ (this is possible as $\mathcal{L}$ is ample, see
Properties, Proposition \href{https://stacks.math.columbia.edu/tag/01Q3}{proposition characterize ample (Tag 01Q3)}).
Then $s^{\deg(t)} \otimes t^{\deg(s)}$
is a nonzero global section of $\mathcal{O}_Z$
(because $Z$ is integral) and hence
a unit (Lemma \href{https://stacks.math.columbia.edu/tag/0BUG}{lemma proper geometrically reduced global sections (Tag 0BUG)}).
This implies that $t$ is a trivializing section
of a positive power of $\mathcal{L}$.
Thus the function $n \mapsto \dim_k \Gamma(X, \mathcal{L}^{\otimes n})$
is bounded on an infinite set of positive integers which contradicts
asymptotic Riemann-Roch
(Proposition \href{https://stacks.math.columbia.edu/tag/0BJ8}{proposition asymptotic riemann roch (Tag 0BJ8)})
since $\dim(Z) > 0$.
\end{proof}


\begin{definition}
\label{divisors-definition-zero-scheme-s}
Let $X$ be a scheme. Let $\mathcal{L}$ be an invertible sheaf.
Let $s \in \Gamma(X, \mathcal{L})$ be a global section.
The {\it zero scheme} of $s$ is the closed subscheme $Z(s) \subset X$
defined by the quasi-coherent sheaf of ideals
$\mathcal{I} \subset \mathcal{O}_X$ which is the image of the
map $s : \mathcal{L}^{\otimes -1} \to \mathcal{O}_X$.
\end{definition}

\section*{Primary problem: Tag 0FD8}
\begin{lemma}[Enriques-Severi-Zariski]
\label{lemma-vanishin-h1-negative}
Let $k$ be a field. Let $X$ be a proper scheme over $k$. Let $\mathcal{L}$
be an ample invertible $\mathcal{O}_X$-module. Let $\mathcal{F}$ be a
coherent $\mathcal{O}_X$-module. Assume that
for $x \in X$ closed we have $\text{depth}(\mathcal{F}_x) \geq 2$.
Then
$H^1(X, \mathcal{F} \otimes_{\mathcal{O}_X} \mathcal{L}^{\otimes m}) = 0$
for $m \ll 0$.
\end{lemma}

\begin{proof}
Choose a closed immersion $i : X \to \mathbf{P}^n_k$ such that
$i^*\mathcal{O}(1) \cong \mathcal{L}^{\otimes e}$ for some $e > 0$
(see Morphisms, Lemma
\href{https://stacks.math.columbia.edu/tag/01VT}{lemma finite type over affine ample very ample (Tag 01VT)}).
Then it suffices to prove the lemma for
$$
\mathcal{G} =
i_*(\mathcal{F} \oplus \mathcal{F} \otimes \mathcal{L} \oplus \ldots
\oplus \mathcal{F} \otimes \mathcal{L}^{\otimes e - 1})
\quad\text{and}\quad \mathcal{O}(1)
$$
on $\mathbf{P}^n_k$. Namely, we have
$$
H^1(\mathbf{P}^n_k, \mathcal{G}(m)) =
\bigoplus\nolimits_{j = 0, \ldots, e - 1}
H^1(X, \mathcal{F} \otimes \mathcal{L}^{\otimes j + me})
$$
by
Cohomology of Schemes, Lemma \href{https://stacks.math.columbia.edu/tag/089W}{lemma relative affine cohomology (Tag 089W)}.
Also, if $y \in \mathbf{P}^n_k$ is a closed point then
$\text{depth}(\mathcal{G}_y) = \infty$ if $y \not \in i(X)$
and $\text{depth}(\mathcal{G}_y) = \text{depth}(\mathcal{F}_x)$
if $y = i(x)$ because in this case
$\mathcal{G}_y \cong \mathcal{F}_x^{\oplus e}$ as a module over
$\mathcal{O}_{\mathbf{P}^n_k, x}$ and we can use for example
Algebra, Lemma \href{https://stacks.math.columbia.edu/tag/0AUK}{lemma depth goes down finite (Tag 0AUK)}
to get the equality.

\medskip\noindent
Assume $X = \mathbf{P}^n_k$ and $\mathcal{L} = \mathcal{O}(1)$ and
$k$ is infinite. Choose $s \in H^0(\mathbf{P}^1_k, \mathcal{O}(1))$
which determines an exact sequence
$$
0 \to \mathcal{F}(-1) \xrightarrow{s} \mathcal{F} \to \mathcal{G} \to 0
$$
as in Lemma \ref{lemma-exact-sequence-induction}. Since the map
$\mathcal{F}(-1) \to \mathcal{F}$ is affine locally given by
multiplying by a nonzerodivisor on $\mathcal{F}$
we see that for $x \in \mathbf{P}^n_k$ closed we have
$\text{depth}(\mathcal{G}_x) \geq 1$, see
Algebra, Lemma \href{https://stacks.math.columbia.edu/tag/090R}{lemma depth drops by one (Tag 090R)}.
Hence by Lemma \ref{lemma-vanishin-h0-negative}
we have $H^0(\mathcal{G}(m)) = 0$ for $m \ll 0$.
Looking at the long exact sequence of cohomology after twisting
(see Remark \href{https://stacks.math.columbia.edu/tag/0EGK}{remark exact sequence induction cohomology (Tag 0EGK)})
we find that the sequence of numbers
$$
\dim H^1(\mathbf{P}^n_k, \mathcal{F}(m))
$$
stabilizes for $m \leq m_0$ for some integer $m_0$.
Let $N$ be the common dimension of these
spaces for $m \leq m_0$. We have to show $N = 0$.

\medskip\noindent
For $d > 0$ and $m \leq m_0$ consider the bilinear map
$$
H^0(\mathbf{P}^n_k, \mathcal{O}(d)) \times
H^1(\mathbf{P}^n_k, \mathcal{F}(m - d))
\longrightarrow
H^1(\mathbf{P}^n_k, \mathcal{F}(m))
$$
By linear algebra, there is a codimension $\leq N^2$ subspace
$V_m \subset H^0(\mathbf{P}^n_k, \mathcal{O}(d))$ such that
multiplication by $s' \in V_m$ annihilates
$H^1(\mathbf{P}^n_k, \mathcal{F}(m - d))$.
Observe that for $m' < m \leq m_0$ the diagram
$$
\xymatrix{
H^0(\mathbf{P}^n_k, \mathcal{O}(d)) \times
H^1(\mathbf{P}^n_k, \mathcal{F}(m' - d)) \ar[r] \ar[d]^{1 \times s^{m' - m}} &
H^1(\mathbf{P}^n_k, \mathcal{F}(m')) \ar[d]^{s^{m' - m}}\\
H^0(\mathbf{P}^n_k, \mathcal{O}(d)) \times
H^1(\mathbf{P}^n_k, \mathcal{F}(m - d)) \ar[r] &
H^1(\mathbf{P}^n_k, \mathcal{F}(m))
}
$$
commutes with isomorphisms going vertically. Thus $V_m = V$ is
independent of $m \leq m_0$. For $x \in \text{Ass}(\mathcal{F})$
set $Z = \overline{\{x\}}$. For $d$ large enough the linear map
$$
H^0(\mathbf{P}^n_k, \mathcal{O}(d)) \to H^0(Z, \mathcal{O}(d)|_Z)
$$
has rank $> N^2$ because $\dim(Z) \geq 1$ (for example this follows
from asymptotic Riemann-Roch and ampleness $\mathcal{O}(1)$; details
omitted). Hence we can find $s' \in V$ such that $s'$ does not vanish
in any associated point of $\mathcal{F}$ (use that the set
of associated points is finite). Then we obtain
$$
0 \to \mathcal{F}(-d) \xrightarrow{s'} \mathcal{F} \to \mathcal{G}' \to 0
$$
and as before we conclude as before that multiplication by $s'$
on $H^1(\mathbf{P}^n_k, \mathcal{F}(m - d))$ is injective
for $m \ll 0$. This contradicts the choice of $s'$ unless
$N = 0$ as desired.

\medskip\noindent
We still have to treat the case where $k$ is finite.
In this case let $K/k$ be any infinite algebraic field extension.
Denote $\mathcal{F}_K$ and $\mathcal{L}_K$ the pullbacks
of $\mathcal{F}$ and $\mathcal{L}$ to $X_K = \Spec(K) \times_{\Spec(k)} X$.
We have
$$
H^1(X_K, \mathcal{F}_K \otimes \mathcal{L}_K^{\otimes m}) =
H^1(X, \mathcal{F} \otimes \mathcal{L}^{\otimes m}) \otimes_k K
$$
by Cohomology of Schemes, Lemma
\href{https://stacks.math.columbia.edu/tag/02KH}{lemma flat base change cohomology (Tag 02KH)}.
On the other hand, a closed point $x_K$ of $X_K$ maps to a closed point
$x$ of $X$ because $K/k$ is an algebraic extension. The
ring map $\mathcal{O}_{X, x} \to \mathcal{O}_{X_K, x_K}$
is flat (Lemma \href{https://stacks.math.columbia.edu/tag/0C4Y}{lemma change fields flat (Tag 0C4Y)}). Hence we have
$$
\text{depth}(\mathcal{F}_{x_K}) =
\text{depth}(\mathcal{F}_x \otimes_{\mathcal{O}_{X, x}}
\mathcal{O}_{X_K, x_K}) \geq
\text{depth}(\mathcal{F}_x)
$$
by Algebra, Lemma \href{https://stacks.math.columbia.edu/tag/0338}{lemma apply grothendieck module (Tag 0338)}
(in fact equality holds here but we don't need it).
Therefore the result over $k$ follows
from the result over the infinite field $K$ and the proof is complete.
\end{proof}
\end{document}
